\input{stock/en-tete_v5.tex}

\newcounter{chapitre}
\setcounter{chapitre}{0}

\title{\Large Chapitre 0 \newline \Huge titre}

\begin{document}
\input{stock/commands.tex}
\maketitle

On note $u_n$ le nombre d’appels engendrés par un appel \code{puissance(x, n)}, sans tenir compte de cet appel lui-même. 

On a, et on va le démontrer~:
$$ \Aboxed{u_n = \lfloor \log_2(n) \rfloor + 1}$$

On rappelle l’algorithme décrit dans l’énoncé~:

\begin{lstLNat}
    fonction puissance(x : TYPE, n : entier long) TYPE
    debut
        si estNul(n) alors 
            retourner elementNeutre()
        sinon si estpair(n) alors
            retourner carre(puissance(x,n>>1))
        sinon
            retourner produit(x , carre(puissance(x,n>>1)))
    fin
\end{lstLNat}


Étant donné un tel algorithme, on a~:
$$ u_n = \begin{cases*}
    0 & si $n = 0$ \\
    u_{\left\lfloor\frac{n}{2}\right\rfloor} + 1 & sinon\\
\end{cases*}$$

\begin{lemme}{1}{}
    Pour tout $n \in \N$, on a~:
    $$2^{\lfloor\log_2(n)\rfloor} \leq n \leq 2^{\lfloor\log_2(n)\rfloor + 1} - 1$$
\end{lemme}

\begin{demonstration}
    Par définition de la partie entière, on a~:
    \begin{align*}
        &\lfloor\log_2(n)\rfloor \leq \log_2(n) < \lfloor\log_2(n)\rfloor + 1\\
        \text{donc} \qquad & 2^{\lfloor\log_2(n)\rfloor} \leq n < 2^{\lfloor\log_2(n)\rfloor + 1} \qquad &\text{car $n \mapsto 2^n$ est strictement croissante}\\
        \text{donc} \qquad & 2^{\lfloor\log_2(n)\rfloor} \leq n \leq 2^{\lfloor\log_2(n)\rfloor + 1} - 1
    \end{align*}
    D'où le lemme 1.
\end{demonstration}

\begin{lemme}{2}{}
    Pour tout $k \in \N$  et tout entier $n$ tel que $2^k \leq n \leq 2^{k+1}-1$, on a~:
    $$u_n = k + 1$$
\end{lemme}

\begin{demonstration}
    Par récurrence sur $k \in \N$, montrons que~:
    $$\forall n \in \N,\, 2^k \leq n \leq 2^{k+1} - 1 \implies u_n = k + 1$$
    \begin{itemize}
        \item Pour $k = 0$, seul $n = 1$ vérifie $2^k = 1 \leq n \leq 1 = 2^{k+1}-1$. Pour ce $n$, on a bien $u_n = 1$ en raison de l'appel \code{puissance(x,0)} engendré.
        \item Supposons maintenant pour $k \in \N$ fixé avoir~:
        $$\forall n \in \N,\, 2^k \leq n \leq 2^{k+1} - 1 \implies u_n = k + 1$$
        Soit $n \in \N$ tel que $2^{k+1} \leq n \leq 2^{k+2}-1$. Montrons que $u_n = k+2$. On a~:
        \begin{align*}
            2^k &\leq \frac{n}{2} \leq \frac{1}{2}\sum_{i=0}^{k+1}2^i = \frac{1}{2} +  \sum_{i=0}^k 2^i\\
            \text{donc}\qquad 2^k &\leq \left\lfloor \frac{n}{2} \right\rfloor \leq \sum_{i=0}^k 2^i = 2^{k+1} -1
        \end{align*}
        Par hypothèse de récurrence, on a $u_{\left\lfloor \frac{n}{2} \right\rfloor}= k + 1$, et puisqu'on a $u_n = u_{\left\lfloor \frac{n}{2}\right\rfloor} + 1$~:
        $$u_n = k +2$$
        Ce qui montre bien l'hypothèse de récurrence au rang suivant.
    \end{itemize}
    D'où le lemme 2 par principe de récurrence.
\end{demonstration}

Soit finalement $n \in \N$. D'après le lemme 1,
$$2^{\lfloor\log_2(n)\rfloor} \leq n \leq 2^{\lfloor\log_2(n)\rfloor + 1} - 1$$
En appliquant le lemme 2 pour $k = \lfloor\log_2(n)\rfloor$ et $n = n$, on obtient le résultat~:
$$u_n = \lfloor\log_2(n)\rfloor + 1$$

\end{document}